%! Piecewise Linear Learning Functions — Unit 8, Lesson 8.1 — Warm-Up
\documentclass[10pt]{article}
\usepackage{linalg-article}
\usepackage{linalg-boxes}
\newcommand{\TallMath}[1]{$\displaystyle #1\rule[-1.4em]{0pt}{3.2em}$}
\newcommand{\relu}{\mathrm{ReLU}}

\begin{document}

\pageheader{Unit 8, Lesson 8.1}{Warm-Up}
\namedateperiod

\vspace{0.12in}

\begin{notesbox}{Getting Ready: The Fold We Built Last Lesson}
\small Last lesson the \textbf{ReLU} $=\max(0,x)$ gave us one \emph{fold}. Today we shift and add folds to
build functions that bend. Warm up on the three pieces of that idea.

\vspace{4pt}
\textbf{1. The ReLU makes one fold (Lesson 8.0).} Evaluate:
\[
  \relu(3)=\blank{0.9cm},\qquad \relu(-2)=\blank{0.9cm},\qquad \relu(0)=\blank{0.9cm}.
\]
As a graph, one ReLU has exactly \blank{1.2cm} fold (corner), at $x=0$.

\vspace{4pt}
\textbf{2. Shifting the input moves the fold.} Consider $\relu(x-2)$. Evaluate it at two inputs:
\[
  \relu(1-2)=\relu(-1)=\blank{0.9cm},\qquad \relu(4-2)=\relu(2)=\blank{0.9cm}.
\]
This function is flat until $x=\blank{0.9cm}$, then rises --- its fold has \emph{slid} to the right.

\vspace{4pt}
\textbf{3. Adding ReLUs changes the slope.} Compute $\relu(x)+\relu(x)$ at $x=3$ and at $x=-1$:
\[
  \relu(3)+\relu(3)=\blank{1.0cm},\qquad \relu(-1)+\relu(-1)=\blank{1.0cm}.
\]
For $x>0$ this equals $2x$: adding two copies \emph{doubled} the slope.
\end{notesbox}

\end{document}
